\begin{comment}
% ===============================
% Slide 1: Impacto fiscal del SGP
% ===============================
\begin{frame}{Sistema General de Participaciones (SGP)}
\begin{itemize}
  \item El SGP atado a los ingresos corrientes de la Nación genera un \textbf{automatismo de gasto}: mayores ingresos ⇒ mayores transferencias.
  \item Con el Acto Legislativo, el SGP pasará de \textbf{4,0\% del PIB en 2024} a \textbf{6,3\% en 2038}.
  \item Una reforma tributaria orientada a consolidación fiscal pierde eficacia, pues parte del nuevo recaudo se destina al SGP.
\end{itemize}
\end{frame}

% ===============================
% Slide 2: Traslado de competencias
% ===============================
\begin{frame}{Ley de Competencias: Traslado de funciones}
\begin{itemize}
  \item Para contener el impacto fiscal del Acto Legislativo, se requiere que las ET asuman nuevas responsabilidades.
  \item Posibles competencias a transferir:
  \begin{itemize}
    \item Inversión pública (2,0\% del PIB).
    \item Salud (2,4\% del PIB).
    \item SENA e ICBF (0,2\% del PIB).
    \item Universidades públicas (0,4\% del PIB).
  \end{itemize}
  \item Transferencia fiscalmente viable solo si va acompañada de transferencia de competencias equivalentes.
\end{itemize}
\end{frame}
\end{comment}


% ===============================
% Slide 3: Sobretasas locales
% ===============================
\begin{frame}{Autonomía Tributaria Territorial}
\begin{itemize}
  \item Se propone dotar a las ET de la capacidad para establecer sobretasas sobre:
  \begin{itemize}
    \item El Impuesto al Valor Agregado (IVA).  
    \item El Impuesto a la Renta.  
  \end{itemize}
  \item Prerequisito: Fondo de Compensación.
%\end{itemize}

%\begin{itemize}
%  \item Avanzar en \textbf{autonomía tributaria territorial}.
  \item Propuesta en IVA:
  \begin{itemize}
    \item Reducir tarifa general del IVA (19\% → 16\%).
    \item Eliminar ICA
    \item Permitir a distritos y departamentos fijar una sobretasa local.
  \end{itemize}
%  \item Incentivos: los territorios se benefician directamente del recaudo adicional generado por formalización y control.
\end{itemize}
\end{frame}

% ===============================
% Slide 3: Sobretasas locales
% ===============================
\begin{frame}{Autonomía Tributaria Territorial}
\begin{itemize}
%  \item Avanzar en \textbf{autonomía tributaria territorial}.
  \item Propuesta en renta jurídicas:
  \begin{itemize}
    \item Reducir tarifa general de renta jurídica en hasta 9 puntos 
    \item Permitir que las entidades territoriales puedan poner una sobretasa de hasta 9 puntos 
    \item Eliminar rentas de destinación específica del antiguo CREE
    \item Reasignar competencias del antiguo CREE 
    \begin{itemize}
        \item 4,4 Salud
        \item 1,4 SENA
        \item 2,2 ICBF 
        \item 0,4 Infancia
        \item 0,6 Universidades
    \end{itemize}
    %\item Reducir tarifa general del IVA (19\% → 16\%).
    %\item Permitir a distritos y departamentos fijar \textbf{sobretasas locales}.
  \end{itemize}
%  \item Incentivos: los territorios se benefician directamente del recaudo adicional generado por formalización y control.
\end{itemize}
\end{frame}


\begin{comment}
% ===============================
% Slide 4: Fórmula del SGP
% ===============================
\begin{frame}{Revisar la fórmula del SGP}
\begin{itemize}
  \item Problema: definición del SGP como porcentaje de los ingresos corrientes de la Nación.
  \item Alternativas:
  \begin{enumerate}
    \item Redefinirlo como porcentaje del PIB (requiere reforma constitucional).
    \item Excluir ingresos adicionales de reformas tributarias, clasificándolos como \textbf{rentas de destinación específica} para el servicio de la deuda.
  \end{enumerate}
  \item Objetivo: evitar que las reformas tributarias pierdan eficacia y que se repita la “reforma cada dos años”.
\end{itemize}
\end{frame}
\end{comment}